\begin{table}[t]
    \centering
    \resizebox{\textwidth}{!}{%
    \begin{tabular}{@{}lllp{6.4cm}@{}}
        \toprule
        \textbf{Setting} & \textbf{Task} & \textbf{Regime} & \textbf{What a skill \emph{could} be} \\
        \midrule
        \arithbten & base-10 3-digit multiplication & elicitation & re-elicit a latent, reliable skill \\
        \vocabes & Spanish$\to$English (24-word glossary) & elicitation & re-elicit words the model already knows \\
        \arithbnine, \arithbeleven & base-9 / base-11 3-digit multiplication & novelty-procedure & scaffold a latent-but-unreliable procedure \\
        \vocabzeth & invented-language (``\zeth'')$\to$English & novelty-information & supply information the model provably cannot have (glossary randomly generated at run time) \\
        \bottomrule
    \end{tabular}
    }
    \caption{The controlled substrate. We dial the regime from \emph{elicitation} (procedure is
    in-distribution, a skill can only re-elicit) to \emph{novelty}. \vocabzeth is the cleanest
    novelty case: its 24 word mappings are randomly generated per run, so they cannot appear in
    any training corpus, yet answers are deterministically checkable.}
    \label{tab:substrate}
\end{table}
